\documentclass[t,12pt]{beamer}

\usepackage{amsmath}
\usepackage{amsfonts}
\usepackage{amssymb}
\usepackage{amsthm}
\usepackage{mathrsfs}
\usepackage{dsfont}
\usepackage[utf8x]{inputenc}
\usepackage[labelfont=bf,justification=justified]{caption}
\usepackage{subcaption}
\usepackage{fixltx2e}
\usepackage{setspace}
\usepackage{arydshln}
\usepackage{xcolor}
\usepackage{tikz}
\usepackage{chngpage}
\usepackage{comment}
\usetikzlibrary{shapes,arrows}
\tikzstyle{block} = [rectangle, draw, fill=blue!20, 
    text width=5em, text centered, rounded corners, minimum height=4em]
\tikzstyle{line} = [draw, -latex']

\usetheme{Emory}
\makeatletter
\setbeamertemplate{footline}
{
  \leavevmode%
  \hbox{%
  \begin{beamercolorbox}[wd=.333333\paperwidth,ht=2.25ex,dp=1ex,center]{author in head/foot}%
    \usebeamerfont{author in head/foot}\insertshortauthor~~\beamer@ifempty{\insertshortinstitute}{}{\insertshortinstitute}
  \end{beamercolorbox}%
  \begin{beamercolorbox}[wd=.333333\paperwidth,ht=2.25ex,dp=1ex,center]{title in head/foot}%
    \usebeamerfont{title in head/foot}\insertshorttitle
  \end{beamercolorbox}%
  \begin{beamercolorbox}[wd=.333333\paperwidth,ht=2.25ex,dp=1ex,right]{date in head/foot}%
    \usebeamerfont{date in head/foot}\insertshortdate{}\hspace*{2em}
    \insertframenumber{} / \inserttotalframenumber\hspace*{2ex} 
  \end{beamercolorbox}}%
  \vskip0pt%
}
\makeatother

\setbeamerfont{frametitle}{size=\large}
\setbeamertemplate{navigation symbols}{}

\titlegraphic{\includegraphics[width=0.3\textwidth]{~/Dropbox/Emory.png}}

\begin{document}

\title[The ARMA framework]{The ARMA framework}
\author{\vspace{0.35in} \\ {\fontfamily{ppl}\selectfont William Mann} \vspace{-0.35in}}

\date{}
\maketitle

\begin{frame}[c]{Outline of next two weeks}
We will discuss several types of time-series processes:
\begin{itemize}
\smallskip
\item AR(p): \emph{Autoregressive} process of order $p$.
\item MA(q): \emph{Moving-average} process of order $q$.
\item ARMA(p,q): Combination of $AR(p)$ and $MA(q)$.
\end{itemize}
\bigskip
These can be combined to approximate any stationary process.
\\ \bigskip
This will us an alternative way to build forecasts, compared to the exponential-smoothing approach from before, 
as well as a sound framework for any other analysis we might want to pursue.
\end{frame}

\begin{frame}[c]
\vfill
\centerline{AR and MA processes, and the Wold theorem}
\vfill
\end{frame}

\begin{frame}[c]{AR(p) process}
Last week, we saw the idea of an AR(1) process:
$$
y_t \quad = \quad \mu \quad + \quad \phi y_{t-1} \quad + \quad u_t
$$
Every date $t$, the new value of $y_t$ is affected by today's shock, 
\\ but also puts some weight $\phi$ on yesterday's $y_{t-1}$.
\\ \bigskip
We also saw the more general version, an AR(p) process:
$$
y_t \,\, = \,\, \mu \,\, + \,\, \phi_1 y_{t-1} \,\, + \,\, \phi_2 y_{t-2} \,\, + \,\, \dots \,\, + \,\, \phi_p y_{t-p} \,\, + \,\, u_t
$$
Recall that we can estimate the $\phi$ numbers by regression.
\\ 
AR(p) is one of the building blocks of the ARIMA framework.
\end{frame}

\begin{frame}[c]{MA(q) process}
The other major building block of time series models is the moving-average or MA process. 
An MA(1) looks like this:
$$
y_t \quad = \quad \mu \quad + \quad u_t \quad + \quad \theta u_{t-1}
$$
Like an AR process, there is persistence built into this process.
\\ 
Differently from AR, the persistence is through the shocks $u_t$.
\\ \medskip
In an MA(1) process, each shock affects today, and to a lesser extent tomorrow, then disappears from the process and no longer matters.
\\ \medskip
By contrast, in an AR process, each shock $u_t$ feeds into $y_t$, 
\\ so it has a long-lived and slowly-decaying effect on future $y$.
\\ \medskip
Again, there can be more than one lag. A general MA(q) process is:
$$
y_t \,\, = \,\, \mu \,\, + \,\, u_t \,\, + \,\, \theta_1 u_{t-1} \,\, + \,\, \theta_2 u_{t-2} + \,\, \dots \,\, + \,\, \theta_q u_{t-q}
$$
\end{frame}

\begin{frame}[c]{Wold representation theorem}
AR and MA processes are useful building blocks, due to the following:
\\ \medskip
\emph{ \textbf{Theorem:} Any covariance-stationary process can be written as a sum of AR(p) and MA(q) processes, plus a non-random constant.}
\\ \medskip
That is, regardless of the true nature behind any process, 
\\ we can always \emph{model} it accurately with something that looks like:
\begin{equation*}
y_t \quad = \quad \mu \, + \, 
\underbrace{ \phi_1 y_{t-1} + \dots + \phi_p y_{t-p} }_{\text{AR(p) terms}}
\, + \,
\underbrace{ u_t + \theta_1 u_{t-1} + \dots + \theta_q u_{t-q} }_{\text{MA(q) terms}}
\end{equation*}
A process that combines AR(p) and MA(q) is called ARMA(p,q).
\vfill
{\scriptsize
The theorem above is slightly paraphrased.
The true theorem only says every stationary process can be written as an MA.
But there is also a ``duality'' between AR and MA, meaning that each one can be converted into the other.
So the choice between them is really just for convenience.
}
\end{frame}

\begin{frame}[c]
\vfill
\centerline{ACF and PACF patterns of AR and MA processes}
\vfill
\end{frame}

\begin{frame}[c]{Autocovariances, autocorrelations, partial autocorrelations}
For a stationary process, $Cov(y_t,y_s)$ depends only on $t-s$.
\begin{itemize}
\bigskip
\item
One way to say this is that $Cov(y_t,y_s) = \gamma_{s-t}$ for some numbers $\gamma_0,\gamma_1,\gamma_2,...$.
The $\gamma$ are called the \textbf{autocovariances} of $y$.
\medskip
\item The \textbf{autocorrelations} of $y$ are $\tau_{s-t} = \gamma_{s-t} / \gamma_0$.
\\
Each $\tau_s$ is also the coefficient from regressing $y_t$ on $y_{t-s}$.
\medskip
\item The \textbf{partial autocorrelations} $\beta_{s-t}$ of $y$ are the coefficients from regressing $y_t$ on $y_{t-s}$ while controlling for the in-between $y$.
\\ There is a formula for them too, but it's more complicted.
\end{itemize}
\bigskip
Next slide: These numbers help us understand stationary processes.
\end{frame}

\begin{frame}[c]{ACF and PACF patterns in AR and MA processes}
One way to understand a stationary process is through its \emph{autocorrelation function} (ACF) or \emph{partial autocorrelation fn} (PACF).
\\ \medskip
These are plots with lags of time on the horizontal axis, 
\\ and autocorrelations or partial autocorrelations on the vertical.
\\ \medskip
These figures help us judge the AR and MA patterns in a process, 
\\ because we know some theoretical benchmarks.
Examples:
\begin{itemize}
\smallskip
\item \textbf{AR(1)}: $ACF(h) = \phi^h$: a \emph{decaying} pattern; 
\\ $PACF(h) = \phi$ if $h=1$, 0 otherwise: a \emph{cutoff} pattern.
\smallskip
\item \textbf{MA(1)}: $ACF(h) = \theta$ if $h=1$, 0 otherwise: a cutoff pattern;
\\ $PACF(h) = \frac{\theta^h(-1)^{h+1}}{1+\theta^2+\dots+\theta^{2h}}$: an \emph{alternating} pattern.
\end{itemize}
\smallskip
See examples in Python notebook posted online.
\end{frame}

\begin{frame}[c]
\vfill
\centerline{Estimating ARMA processes}
\vfill
\end{frame}

\begin{frame}[c]{Estimating an AR($p$) process}
Suppose you have a series of data from a stationary process,
\\ and you think the process is described well by an AR($p$):
$$
y_t \,\, = \,\, \mu \,\, + \,\, \phi_1 y_{t-1} \,\, + \,\, \phi_2 y_{t-2} \,\, + \,\, \dots \,\, + \,\, \phi_p y_{t-p} \,\, + \,\, u_t
$$
How would you go about estimating the numbers $\phi_1$, $\phi_2$, ... $\phi_p$?
\\ \medskip
We've seen the answer before: Just run a regression. 
\\ See examples on Canvas.
\end{frame}

\begin{frame}[c]{Estimating an MA($q$) process}
Suppose you have a series of data from a stationary process,
\\ and you think the process is described well by an MA($q$):
$$
y_t \,\, = \,\, \mu \,\, + \,\, u_t \,\, + \,\, \theta_1 u_{t-1} \,\, + \,\, \theta_2 u_{t-2} + \,\, \dots \,\, + \,\, \theta_q u_{t-q}
$$
How would you go about estimating the numbers $\theta_1$, $\theta_2$, ... $\theta_p$?
\\ \bigskip
Unfortunately, you \textbf{cannot} run any simple regression for this.
\\ You can't actually observe the ``true'' shocks $u_t$ that appear above,
\\ and you can't estimate them until you have an estimate of $\theta$.
\\ \bigskip
We use a different estimation method called \textbf{maximum likelihood}.
\end{frame}

\begin{frame}[c]{Maximum likelihood estimation}
Maximum likelihood is an important approach in many settings. 
\\ (We will see it again with volatility models in a few weeks.)
\\ \bigskip The basic idea is this:
\begin{itemize}
\medskip
\item Make an initial guess of all the model parameters:
\\ $\mu$, $\sigma^2$, and any AR and MA terms $\phi_1,\phi_2,...,\theta_1,\theta_2,...$.
\smallskip
\item Calculate the \textbf{likelihood} (probability) that you would observe the data you did, if the parameters were equal to your guess.
\smallskip
\item Try different guesses until that likelihood is as high as possible.
\end{itemize}
\bigskip
The parameter values that generate the highest likelihood of observing the true data are the \textbf{maximum likelihood estimate}.
\end{frame}

\begin{frame}[c]{The likelihood function}
\begin{itemize}
\item To implement MLE, you need to know the probability described in the prior slide, which is called the \textbf{likelihood function}. 
\medskip
\item For any standard ARMA process, the function is well-known.
\\ You just have a computer find parameters to maximize it.
\\ More often, you actually maximize the \emph{log} of the likelihood.
\medskip
\item Example: Here's the log likelihood function for an AR(1).
\begin{align*}
\mathscr{L} &= - \frac{T}{2}\ln(\sigma^2) - \frac{T}{2}\ln(2\phi) - \frac{1}{2}\ln\left(\frac{\sigma^2}{1-\phi^2}\right) 
\\ & \quad - \frac{1-\phi^2}{2\sigma^2}\hat{u}_1^2 - \frac{1}{2\sigma^2}\sum_{t=2}^T \hat{u}_t^2
\end{align*}
\end{itemize}
\end{frame}

\begin{frame}[c]
\vfill
\centerline{After estimation: Measures of model quality}
\vfill
\end{frame}

\begin{frame}[c]{Information criteria}
An \emph{information criterion} compares models with each other:
\begin{itemize}
\medskip
\item First estimate the model, and measure its \emph{fit} in the data. 
\begin{itemize}
\item For a linear regression, fit is measured by $R^2$. 
\item When applying MLE, we use the likelihood function, evaluated at our best parameter estimates. 
Higher value $=$ better fit.
\item Most software reports the \emph{log} of the likelihood function, which is typically negative, so a \emph{less-negative} number is a better fit.
\end{itemize}
\medskip
\item Build in some penalty to avoid \emph{overfitting}.
\begin{itemize}
\item The penalty should increase with the number of model terms.
\item In linear regression, adjusted $R^2$ is an example of this.
\end{itemize}
\end{itemize}
\medskip
The two standard information criteria are AIC and BIC.
\\ See next two slides and examples in Python notebook.
\end{frame}

\begin{frame}[c]{Akaike information criterion (AIC)}
The AIC measure of a model is given by
$$
AIC \quad = \quad 2 \times k \quad - \quad 2 \times \mathscr{L}
$$
\begin{itemize}
\vspace{-0.2in}
\item $k$ is the number of parameters that were estimated.
\begin{itemize}
\item For an ARMA model, this is the number of AR and MA terms, plus 1 for $\sigma^2$, plus 1 if there was a constant term $\mu$.
\item Example: ARMA(1,1) with a constant term would have $k=4$.
\end{itemize}
\medskip
\item $\mathscr{L}$ is the maximum \emph{log} likelihood that we can achieve.
\begin{itemize}
\item This is part of the output from your MLE process.
\end{itemize}
\medskip
\item You want to find the \textbf{smallest} possible AIC.
\begin{itemize}
\item Intuition: Higher $\mathscr{L}$ is good. Higher $k$ is bad.
\end{itemize}
\end{itemize}
\end{frame}

\begin{frame}[c]{Bayesian information criterion (BIC, SIC, SBC, or SBIC)}
The BIC measure of a model is given by
$$
BIC \quad = \quad \ln(N) \times k \quad - \quad 2 \times \mathscr{L}
$$
\begin{itemize}
\vspace{-0.2in}
\item $N$ is the number of observations in the dataset,
\\ everything else is as in the prior slide.
\medskip
\item Intuition: Just a different way of penalizing $k$.
\medskip 
\item There are deep statistical arguments for why the AIC and BIC look the way they do, but this is beyond the scope of class.
\medskip
\item Both are popular and should yield similar conclusions.
\end{itemize}
\end{frame}

\begin{frame}[c]{How to use information criteria}
It's common to search over models to find the smallest AIC or BIC.
\begin{itemize}
\bigskip
\item Combine common sense and judgment with this process.
\begin{itemize}
\item Many times, two different models will yield very similar AIC and BIC numbers. It may be that one model scores marginally better, but the other is more economically reasonable.
\end{itemize}
\medskip
\item AIC and BIC only give \emph{relative} comparisons. 
\\ There is no sense of a ``good'' or ``bad'' value to judge by.
\begin{itemize}
\item For \emph{overall} assessment of a model, you should instead check whether there are any remaining features of the data that seem important to capture.
See next section.
\end{itemize}
\end{itemize}
\end{frame}

\begin{frame}[c]
\vfill
\centerline{After estimation: Check for remaining time-series patterns}
\vfill
\end{frame}

\begin{frame}[c]{Checking for time-series patterns}
We are done with model selection, once we are satisfied that we have captured all the important time-series features of the data.
\\ \bigskip
To judge this, we examine the \emph{residuals} from the fitted models.
\begin{itemize}
\medskip
\item Visually, look at ACF and PACF graphs of the residuals.
\medskip
\item Formally, we can apply statistical tests.
\\ The most common choice is the Ljung-Box test.
\end{itemize}
\end{frame}

\begin{frame}[c]{Ljung-Box test}
The Ljung-Box test statistic is defined as 
$$
Q = n(n+2) \sum_{k=1}^h \frac{\hat{\rho}_k^2}{n-k}
$$
where $\hat{\rho}_k$ is the sample autocorrelation up to lag $k$, and $h$ is some specified number of lags we want to include.
\\ \medskip
The null hypothesis is that the data are independently distributed. 
\\ Under this hypothesis, $Q$ follows a standard $\chi^2$ distribution.
\\ \medskip
Hence, compare $Q$ with critical values of this distribution.
\\ \medskip
The test will tell us if there is \emph{some} remaining time series feature left in the data (up to lag $h$), but will not tell us exactly what it is.
\end{frame}

\begin{frame}[c]{Common practice}
\begin{itemize}
\item The Ljung-Box test statistic and p-values are a standard part of the output from ARMA estimation (usually with $h=1$).
\medskip
\item A common approach is to try different models until L-B test fails to reject the null that the residuals have no persistence.
\medskip
\item Again, keep in mind the role of common sense and judgment. 
\\
If you already have a sensible model, 
it may not be worth adding extra complications just to lower the L-B statistic marginally.
\end{itemize}
\medskip
\begin{itemize}
\item 
{\footnotesize
\textbf{Critique of common practice:}
When there are AR terms in the model, it is actually \emph{incorrect} to apply the L-B test to the residuals of the model.
The correct test is a modified version called the \emph{Breusch-Godfrey} test, which leads to lower p-values.
But you will rarely (if ever) see this in practice.
}
\end{itemize}
\end{frame}

\begin{frame}[c]
\vfill
\centerline{ARMA model selection: The Box-Jenkins approach}
\vfill
\end{frame}

\begin{frame}[c]{The Box-Jenkins approach}
Combining all our pieces so far, we can map out an ARMA-based framework for modeling and forecasting a stationary data series:
\begin{itemize}
\smallskip
\item Examine the series, especially its ACF and PACF graphs, 
\\ to determine a few choices of $p$ and $q$ that we should consider.
\item For each pair of $(p,q)$, estimate the parameter values by MLE. 
\item Compare each set of results based on AIC or BIC. 
\\ Check the ACF and PACF of the residuals, and the L-B test, 
\\ for any remaining time-series patterns that should be modeled.
\item Choose a model that scores well and is economically reasonable.
\end{itemize}
\smallskip
If the data are not stationary, you first transform as necessary so that they are, and then follow this procedure. See discussion later.
\end{frame}

\begin{frame}[c]
\vfill
\centerline{After model selection and estimation:}
\centerline{Building forecasts and prediction intervals}
\vfill
\end{frame}

\begin{frame}[c]{One-step forecast and prediction interval for MA(1)}
Let's first think about a simple MA(1) process $y_t = \mu + u_t + \theta u_{t-1}$.
\\ \smallskip
Suppose the data ends at $T$. Let's forecast $y_{T+1} = \mu + u_{T+1} + \theta u_T$.
\begin{itemize}
\smallskip
\item We first fit the model to get estimates of $\mu$, $\theta$, and also all of the $u$ values (using the sample residuals from the model).
\begin{itemize}
\smallskip
\item All of these are just estimates, meaning there is still uncertainty in them -- but we will ignore that and act as though they are the true values. This is not quite ideal, but it is standard.
\end{itemize}
\smallskip
\item As of time $T$, we forecast $y_{T+1}$ with its conditional expectation,
\\ $\mathbb{E}_T[y_{T+1}] = \mu + \theta u_T$ (using the fact that $\mathbb{E}[u_{T+1}]=0$).
\medskip
\item $Var_T(y_{T+1}) = \sigma^2$, since $u_{T+1}$ is the only thing unknown as of $t$.
\\ A 95\% prediction interval is our forecast $\pm 2 \times \sqrt{\sigma^2}$.
\end{itemize}
\end{frame}

\begin{frame}[c]{Two-step forecast and prediction interval for MA(1)}
Now let's forecast $y_{T+2} = \mu + u_{T+2} + \theta u_{T+1}$ as of time $T$.
\\ \bigskip
The forecast is $\mathbb{E}_T[y_{T+2}] = \mu$ since we don't know $u_{T+2}$ or $u_{T+1}$.
\\ \bigskip
Also because we don't know either of these errors, the prediction interval is wider:
$Var_T(y_{T+2}) = Var(u_{T+2} + \theta u_{T+1}) = \sigma^2\times(1+\theta^2)$.
\\ \bigskip
A 95\% confidence interval is our forecast $\pm \sqrt{ \sigma^2\times(1+\theta^2) }$.
\\ \bigskip
In fact this expected value and variance are the long-run values, 
so they will apply to all dates $T+2$ onward.
This reflects the fact that shocks only affect an MA(1) process for two dates before dying out.
\end{frame}

\begin{frame}[c]{Forecasting an AR(1) process}
Now consider an AR(1) process, $y_t = \mu + \phi y_{t-1} + u_t$.
\\ \bigskip
At time $T$ we forecast $y_{T+1} = \mu + \phi y_T + u_t$ as $\mathbb{E}_T[y_{T+1}] = \mu + \phi y_T$.
\\ \bigskip
Farther out, $\mathbb{E}_T[y_{T+2}] = \mathbb{E}_T[\mu + \phi y_{T+1} + u_T] = \mu + \phi \mathbb{E}_T[y_{T+1}]$.
\\ \smallskip And we just calculated $\mathbb{E}_T[y_{T+1}] = \mu + \phi y_T$.
\\ \smallskip We can just substitute that in, which gives us
$$
\mathbb{E}_T[y_{T+2}] = \mu(1+\phi) + \phi^2 y_T
$$
We can repeat this process to forecast as far out as we want.
\\ \medskip
Not surprisingly, the forecasts for an AR(1) will start close to $y_T$, 
\\ and decay slowly to the long-run average of $\frac{\mu}{1-\phi}$.
\end{frame}

\begin{frame}[c]{Prediction interval for an AR(1) process}
Just think carefully about what is known and unknown at time $T$:
\begin{align*}
Var_T(y_{T+1}) &= Var_T(\mu + \phi y_{T} + u_{T+1}) = \sigma^2
\\
Var_T(y_{T+2}) &= Var_T(\mu + \phi y_{T+1} + u_{T+2}) = \sigma^2 \phi^2 + \sigma^2 = \sigma^2(1+\phi^2)
\\
Var_T(y_{T+3}) &= Var_T(\mu + \phi y_{T+2} + u_{T+3}) 
\\
&\quad = \phi^2 [ \sigma^2(1 + \phi^2)] + \sigma^2 = \sigma^2(1+\phi^2+\phi^4)
\end{align*}
...and so forth. The variance grows towards a limit of $\frac{\sigma^2}{1-\phi^2}$.
\end{frame}

\begin{frame}[c]{Forecasts and prediction intervals in general}
The procedure is the same for any ARMA model:
\begin{itemize}
\smallskip
\item Build forecasts by substituting in values of $y$ and $u$ from the end of the dataset into the AR and MA terms of the model. 
\begin{itemize}
\item After reaching the last $y$, substitute forecasts of $y$.
\item After reaching the last $u$, substitute zeros.
\end{itemize}
\smallskip
\item For a prediction interval at horizon $k$:
\begin{itemize}
\item Add up $\phi^{2(k-1)} + \phi^{2(k-2)} ... $, and multiply by $\sigma^2$. 
This is the forecast variance.
Take the square root of that variance to get the standard error. A 95\% CI is the forecast $\pm$ two SEs.
\end{itemize}
\smallskip
\item Of course this should really be done with software, not by hand.
\end{itemize}
\end{frame}

\begin{frame}[c]{Practical notes}
\begin{itemize}
\item All the above formulas for prediction intervals assume that the errors in the model are independent.
They are not quite correct if there are remaining time-series patterns in the residuals.
\begin{itemize}
\item 
This is one reason to check for such patterns. 
\item However, the \emph{forecasts} themselves are not affected by this issue, as they only assume the shocks are zero on average.
\end{itemize}
\medskip
\item \emph{All} the formulas ignore the fact that the model parameters themselves are estimated with uncertainty.
\begin{itemize}
\item This is a limitation, but it is standard. 
\item One result is that prediction intervals are always too narrow.
\end{itemize}
\end{itemize}
\end{frame}

\begin{frame}[c]
\vfill
\centerline{Handling trends and seasonality in the ARMA framework}
\vfill
\end{frame}

\begin{frame}[c]{Comparing ARMA with exponential smoothing methods}
The ARMA framework assumes you have a stationary process.
\begin{itemize}
\item This means that ARMA is best seen as an alternative to \emph{simple} exponential smoothing.
It will fail to capture stochastic trends, or seasonality, for the same reason that SES failed to do so.
\end{itemize}
\medskip
Recall how we extended SES to address these issues:
\begin{itemize}
\item To handle trends, we used ``double'' exponential smoothing.
\item To handle trends \emph{and} seasonality, we introduced ``triple'' exponential smoothing (also called the Holt-Winters method).
\end{itemize}
\medskip
How do we do the same thing in the ARMA framework?
\begin{itemize}
\item To handle trends, we extend to ``ARIMA'' models.
\\ Basically, modeling \emph{differenced} data by ARMA.
\item To handle trends \emph{and} seasonality, we extend to ``SARIMA,''
which features ``seasonal'' AR, MA, and differencing terms.
\end{itemize}
\end{frame}

\begin{frame}[c]{Trending data: ARIMA models}
A strong trend is a source of nonstationary behavior.
\\ \bigskip
In the ARMA framework, trends are typically modeled as a unit root:
\begin{itemize}
\item
Difference the data, and check that this difference is stationary.
\item
Then analyze the differenced data with the ARMA procedure.
\end{itemize}
\bigskip
A process whose difference is ARMA$(p,q)$ is called ARIMA$(p,d,q)$.
\\ ARIMA stands for \textbf{A}uto\textbf{R}egressive \textbf{I}ntegrated \textbf{M}oving \textbf{A}verage.
\begin{itemize}
\item
The $p$ and $q$ parameters have the same meaning as before.
\item
We add $d$ in between them to denote the number of times that the data must be differenced in order for it to be stationary.
\item
In almost every real-world scenario, $d=1$.
\end{itemize}
\end{frame}

\begin{frame}[c]{Seasonal data: SARMA and SARIMA models}
\begin{itemize}
\item
Seasonality leads to regular spikes in the ACF and PACF.
\\ These cannot be modeled with standard AR or MA processes.
\medskip
\item
Instead, we add another dimension to the model:
\\ A seasonal ARIMA model (SARIMA) has the usual ARIMA terms,
plus \emph{seasonal} versions of these terms labeled $(P,D,Q)_s$.
\medskip
\item
Altogether this is written as SARIMA$(p,d,q)\times(P,D,Q)_s$.
\medskip
\item SARIMA is the most general category of ARMA model,
\\ and produces forecasts comparable to Holt-Winters.
\medskip
\item See examples on Canvas of SARIMA estimation and forecasting.
\end{itemize}
\end{frame}

\end{document}
